\documentclass[otchet]{SCWorks}
% Тип обучения (одно из значений):
%    bachelor   - бакалавриат (по умолчанию)
%    spec       - специальность
%    master     - магистратура
% Форма обучения (одно из значений):
%    och        - очное (по умолчанию)
%    zaoch      - заочное
% Тип работы (одно из значений):
%    coursework - курсовая работа (по умолчанию)
%    referat    - реферат
%  * otchet     - универсальный отчет
%  * nirjournal - журнал НИР
%  * digital    - итоговая работа для цифровой кафедры
%    diploma    - дипломная работа
%    pract      - отчет о научно-исследовательской работе
%    autoref    - автореферат выпускной работы
%    assignment - задание на выпускную квалификационную работу
%    review     - отзыв руководителя
%    critique   - рецензия на выпускную работу

% * Добавлены вручную. За вопросами к @mchernigin
\usepackage{preamble}

\begin{document}

% Кафедра (в родительном падеже)
\chair{математической кибернетики и компьютерных наук}

% Тема работы
\title{Двоичное дерево поиска.}

% Курс
\course{3}

% Группа
\group{351}

% Факультет (в родительном падеже) (по умолчанию "факультета КНиИТ")
% \department{факультета КНиИТ}

% Специальность/направление код - наименование
% \napravlenie{02.03.02 "--- Фундаментальная информатика и информационные технологии}
% \napravlenie{02.03.01 "--- Математическое обеспечение и администрирование информационных систем}
% \napravlenie{09.03.01 "--- Информатика и вычислительная техника}
\napravlenie{09.03.04 "--- Программная инженерия}
% \napravlenie{10.05.01 "--- Компьютерная безопасность}

% Для студентки. Для работы студента следующая команда не нужна.
\studenttitle{студентки}

% Фамилия, имя, отчество в родительном падеже
\author{Винник Маргариты Дмитриевны}

% Заведующий кафедрой 
\chtitle{доцент, к.\,ф.-м.\,н.}
\chname{С.\,В.\,Миронов}

% Руководитель ДПП ПП для цифровой кафедры (перекрывает заведующего кафедры)
% \chpretitle{
%     заведующий кафедрой математических основ информатики и олимпиадного\\
%     программирования на базе МАОУ <<Ф"=Т лицей №1>>
% }
% \chtitle{г. Саратов, к.\,ф.-м.\,н., доцент}
% \chname{Кондратова\, Ю.\,Н.}

% Научный руководитель (для реферата преподаватель проверяющий работу)
\satitle{доцент, к.\,ф.-м.\,н.} %должность, степень, звание
\saname{И.\,А.\,Батраева}

% Руководитель практики от организации (руководитель для цифровой кафедры)
\patitle{доцент, к.\,ф.-м.\,н.}
\paname{С.\,В.\,Миронов}

% Руководитель НИР
\nirtitle{доцент, к.\,п.\,н.} % степень, звание
\nirname{В.\,А.\,Векслер}

% Семестр (только для практики, для остальных типов работ не используется)
\term{2}

% Наименование практики (только для практики, для остальных типов работ не
% используется)
\practtype{учебная}

% Продолжительность практики (количество недель) (только для практики, для
% остальных типов работ не используется)
\duration{2}

% Даты начала и окончания практики (только для практики, для остальных типов
% работ не используется)
\practStart{01.07.2024}
\practFinish{13.01.2024}

% Год выполнения отчета
\date{2026}

\maketitle

% Включение нумерации рисунков, формул и таблиц по разделам (по умолчанию -
% нумерация сквозная) (допускается оба вида нумерации)
\secNumbering

\tableofcontents

% Раздел "Обозначения и сокращения". Может отсутствовать в работе
% \abbreviations
% \begin{description}
%     \item ... "--- ...
%     \item ... "--- ...
% \end{description}

% Раздел "Определения". Может отсутствовать в работе
% \definitions

% Раздел "Определения, обозначения и сокращения". Может отсутствовать в работе.
% Если присутствует, то заменяет собой разделы "Обозначения и сокращения" и
% "Определения"
% \defabbr

\intro

Двоичное дерево поиска (англ.\ \emph{binary search tree}, BST) "--- структура
данных для хранения упорядоченного динамического множества. Она поддерживает
поиск, вставку и удаление ключа, а также поиск минимума, максимума,
предшественника и преемника~\cite{kormen}. На деревьях поиска, как правило,
строятся ассоциативные контейнеры: \texttt{std::set} и \texttt{std::map} в
C++, \texttt{TreeSet} в Java.
 
Время работы основных операций пропорционально высоте дерева $h$: для
сбалансированного дерева $h = \Theta(\lg n)$, а для вырожденного дерева,
напоминающего связный список, $h = \Theta(n)$~\cite{kormen}. Поэтому простое
двоичное дерево поиска не даёт гарантий, и на практике используют его
самобалансирующиеся версии: красно-чёрные деревья и АВЛ-деревья. Им будут
посвящены следующие лабораторные работы.
 
\textbf{Цель работы} "--- изучить двоичное дерево поиска, реализовать его на
языке C++ и оценить сложность операций и расход памяти.
 
\textbf{Задачи:}
\begin{enumerate}
  \item изучить определение двоичного дерева поиска, его свойство и способы обхода;
  \item реализовать наивную версию (вставка, поиск, вывод) и проверить её на тестовых задачах;
  \item разобрать готовую реализацию Simple Snippets (удаление, три вида обхода, обход по уровням, высота) и выявить её недостатки;
  \item написать улучшенную реализацию, устранив выявленные недостатки;
  \item оценить сложность операций и расход памяти.
\end{enumerate}

\section{Краткая теория}

\subsection{Определения}
 
\emph{Дерево} "--- непустая коллекция вершин и рёбер, в которой любые две
вершины соединены ровно одним путём~\cite{sedgewick}. \emph{Двоичное
дерево} "--- дерево, в котором у каждой вершины не более двух потомков:
левый и правый. \emph{Двоичное дерево поиска} "--- двоичное дерево, в котором
для любого узла $x$ ключи левого поддерева меньше $x.\mathit{key}$, а ключи
правого поддерева больше~\cite{itmo}.
 
В~\cite{kormen} допускаются равные ключи (в левом поддереве $\le$, в правом
$\ge$), поэтому на рисунке~\ref{fig:bst-example} есть два узла со значением~5.
В нашей реализации все ключи уникальны, поэтому неравенства строгие.
 
\subsection{Представление в памяти}
 
Узел хранит ключ \texttt{key} и указатели \texttt{left\_child} и
\texttt{right\_child}; отсутствующий потомок обозначается
\texttt{nullptr}. В ряде реализаций в узел добавляют указатель на родителя
$p$~\cite{kormen}, но в данной работе информация о родителе не используется.
Всё дерево задаётся указателем на корень \texttt{root}.
 
\subsection{Основные операции}
 
\begin{itemize}
  \item \textbf{Поиск.} Сравниваем искомый ключ с ключом текущего узла и
        переходим в левое или правое поддерево, пока не найдём ключ или не
        дойдём до \texttt{nullptr}. Время $O(h)$.
  \item \textbf{Вставка.} Спускаемся так же, как при поиске, и подвешиваем
        новый узел на месте \texttt{nullptr}. Время $O(h)$.
  \item \textbf{Удаление.} Рассматриваются три случая: у узла нет детей
        (узел просто удаляется); у узла один ребёнок (узел заменяется
        ребёнком); у узла два ребёнка (на его место ставится преемник "---
        минимальный узел правого поддерева). Время $O(h)$.
  \item \textbf{Обходы.} \emph{Прямой} (узел, левое, правое),
        \emph{симметричный} (левое, узел, правое; ключи выводятся по
        возрастанию) и \emph{обратный} (левое, правое, узел). Время
        $\Theta(n)$~\cite{kormen}.
\end{itemize}
 
\subsection{Зависимость времени работы от высоты}
 
Одно и то же множество ключей может быть представлено разными деревьями
поиска (рисунок~\ref{fig:bst-example}): в варианте (а) высота равна~2, в
варианте (б) "--- 4. Форма дерева определяется порядком вставки ключей:
для случайного порядка математическое ожидание высоты равно $O(\lg n)$, а
для отсортированной последовательности дерево превращается в цепочку высотой
$n-1$~\cite{kormen}. Именно поэтому нужны самобалансирующиеся деревья.
 
% Если в шаблоне подпись рисунка выводится как «Рис. 1.», то добавьте в
% преамбулу: \addto\captionsrussian{\renewcommand{\figurename}{Рисунок}}
% \captionsetup[figure]{labelsep=endash}  % получится «Рисунок 1 -- ...»
\begin{figure}[htbp]
  \centering
  \begin{tikzpicture}[
      every node/.style={circle, draw, minimum size=7mm, inner sep=1pt},
      level distance=11mm,
      level 1/.style={sibling distance=24mm},
      level 2/.style={sibling distance=13mm},
      level 3/.style={sibling distance=11mm},
      level 4/.style={sibling distance=9mm}]
    % (а)
    \node (a) at (0,0) {6}
      child {node {5} child {node {2}} child {node {5}}}
      child {node {7} child[missing] child {node {8}}};
    \node[draw=none, minimum size=0pt] at (0,-3.9) {(а)};
    % (б)
    \node (b) at (8,0) {2}
      child[missing]
      child {node {5}
        child[missing]
        child {node {7}
          child {node {6} child {node {5}} child[missing]}
          child {node {8}}}};
    \node[draw=none, minimum size=0pt] at (8.6,-5.2) {(б)};
  \end{tikzpicture}
  \caption{Двоичные деревья поиска с одним и тем же набором ключей (по~\cite{kormen})}
  \label{fig:bst-example}
\end{figure}

\subsection{Время работы обходов}
\label{sec:walk-time}
 
Все три обхода (прямой, симметричный и обратный) выполняются за время,
линейное по числу узлов. Приведём доказательство из~\cite{kormen} для
симметричного обхода. Для двух других оно полностью повторяется: процедуры
отличаются только тем, где относительно двух рекурсивных вызовов
печатается ключ, а на время работы это не влияет.
 
Процедура симметричного обхода:
 
\begin{Verbatim}[fontsize=\small]
INORDER-TREE-WALK(x)
1  if x != NIL
2      INORDER-TREE-WALK(x.left)
3      print x.key
4      INORDER-TREE-WALK(x.right)
\end{Verbatim}
 
\textbf{Теорема.} Если $x$ "--- корень поддерева с $n$ узлами, то время
работы вызова \texttt{INORDER-TREE-WALK}$(x)$ равно $\Theta(n)$.
 
\textbf{Доказательство.} Пусть $T(n)$ "--- время вызова для поддерева с $n$
узлами. Процедура посещает все $n$ узлов, поэтому $T(n) = \Omega(n)$.
Остаётся показать, что $T(n) = O(n)$.
 
Для пустого дерева выполняется только проверка $x \neq \mathrm{NIL}$,
поэтому $T(0) = c$ для некоторой константы $c > 0$. При $n > 0$ процедура
вызывается для левого поддерева с $k$ узлами и для правого с $n-k-1$
узлами, поэтому
\[
  T(n) \le T(k) + T(n-k-1) + d,
\]
где $d > 0$ "--- константа, отражающая время выполнения тела процедуры без
рекурсивных вызовов.
 
Методом подстановки докажем, что $T(n) \le (c+d)n + c$. Для $n = 0$ имеем
$(c+d)\cdot 0 + c = c = T(0)$. Для $n > 0$:
\begin{align*}
  T(n) &\le T(k) + T(n-k-1) + d \\
       &\le \bigl((c+d)k + c\bigr) + \bigl((c+d)(n-k-1) + c\bigr) + d \\
       &= (c+d)(n-1) + 2c + d \\
       &= (c+d)n + c.
\end{align*}
Следовательно, $T(n) = O(n)$, и вместе с оценкой снизу получаем
$T(n) = \Theta(n)$. \hfill $\square$


\section{Наивная реализация}

Для начала рассмотрим так называемую наивную реализацию двоичного дерева
поиска. Определение дерева и его свойства приведены в
разделе~«Краткая теория». Здесь рассматривается вариант без информации о
родителе: узел хранит только ключ \texttt{key} и указатели на потомков
\texttt{left\_child} и \texttt{right\_child}, а само дерево "--- указатель на
корень \texttt{root}.
 
Такая реализация нужна по следующим причинам:
\begin{itemize}
  \item она даёт представление об устройстве \texttt{std::set} в C++,
        \texttt{TreeSet} в Java и других подобных коллекций;
  \item на ней видно, как на практике работает двоичное дерево поиска
        (в теории всё выглядит понятным, а детали проявляются только при
        написании кода);
  \item она показывает недостатки базовой версии и объясняет, зачем нужна
        балансировка;
  \item она служит основой для собственного аналога \texttt{std::set},
        поддерживающего операции, которых нет в стандартном контейнере.
\end{itemize}
 
Реализация является наивной по следующим причинам:
\begin{itemize}
  \item нет балансировки;
  \item высота дерева не контролируется: у структуры нет даже поля,
        хранящего высоту;
  \item все операции работают за $O(h)$, где $h$ "--- высота дерева, а в
        худшем случае $h = O(n)$.
\end{itemize}
 
Версия поддерживает вставку, поиск и вывод дерева. Для проверки решены две
тестовые задачи; их условия приведены в комментариях к коду.

\begin{minted}{c}
#define _CRT_SECURE_NO_WARNINGS
#include <cstdio>
#include <iostream>
#include <string>

using namespace std;
/* 
Правило двоичного дерева поиска:
1) у любой вершины в левом поддереве ключи < ключу в самой вершине
2) у любой вершины в правом поддереве ключи > ключу в самой вершине
*/

/*class BST {
	struct Node
	{
		int key;
		Node *left_child = nullptr, * right_child = nullptr;
		// конструктор. принимает значение ключа и 
		// записывает это значение в поле key
		Node(int k) : key(k) {}
	}
	*root = nullptr;

	bool contains(Node* n, int key) const 
	{
		if (!n)
			return false;
		if (key == n -> key)
			return true;
		if (key < n->key)
			return contains(n->left_child, key);
		else
			return contains(n->right_child, key);
	}
	void insert(Node *n, int key)
	{
		if (!root)
		{
			root = new Node(key);
			return;
		}
		if (key == n ->key)
		{
			return;
		}
		if (key < n->key)
		{
			if (n->left_child)
			{
				insert(n->left_child, key);
			}
			else
			{
				n->left_child = new Node(key);
			}
		}
		else
		{
			if (n->right_child)
			{
				insert(n->right_child, key);
			}
			else
			{
				n->right_child = new Node(key);
			}
		}
	}
public:
	bool contains(int key) const
	{
		return contains(root, key);
	}
	void insert(int key) 
	{
		insert(root, key);
	}
};
*/

class BST {
	struct Node
	{
		int key;
		Node* left_child = nullptr, * right_child = nullptr;
		Node(int k) : key(k) {}
	}
	*root = nullptr;

	bool contains(Node* n, int key) const
	{
		if (!n)
			return false;
		if (key == n->key)
			return true;
		if (key < n->key)
			return contains(n->left_child, key);
		else
			return contains(n->right_child, key);
	}
	// основная функция, от кот. отталкиваемся
	// и модифицируем
	/*void insert(Node*& n, int key)
	{
		if (!n)
		{
			n = new Node(key);
			return;
		}
		if (key == n->key)
		{
			return;
		}
		if (key < n->key)
		{
			insert(n->left_child, key);
		}
		else
		{
			insert(n->right_child, key);
		}
	}
	*/
	// модификация функции вставки под задачу 2
	bool insert(Node*& n, int key)
	{
		if (!n)
		{
			n = new Node(key);
			return true;
		}
		if (key == n->key)
			return false;
		if (key < n->key)
			return insert(n->left_child, key);
		else
			return insert(n->right_child, key);
	}
	// основная функция, от кот.отталкиваемся
	// и модифицируем
	/*
	void print(Node *n) const
	{
		if (!n)
			return;
		print(n->left_child);
		if (n->left_child && n->right_child)
		{
			cout << n->key << "\n";
		}
		//cout << n->key << "\n";
		print(n->right_child);
	}
	*/
	void print(Node* n, int level) const
	{
		if (!n)
			return;
		print(n->left_child, level + 1);
		for (int i = 0; i < level; i++)
		{
			cout << ".";
		}		
		cout << n->key << "\n";
		print(n->right_child, level + 1);
	}
	void destroy(Node *n)
	{
		if (!n)
			return;
		destroy(n->left_child);
		destroy(n->right_child);
		delete n;
	}
public:
	BST() = default;
	BST(const BST&) = delete;
	BST& operator=(const BST&) = delete;

	bool contains(int key) const
	{
		return contains(root, key);
	}
	bool insert(int key)
	{
		return insert(root, key);
	}
	/*
	void insert(int key)
	{
		insert(root, key);
	}
	
	void print() const 
	{
		print(root);
	}
	*/
	void print() const
	{
		print(root, 0);
	}
	~BST() 
	{
		destroy(root);
	}
};

// для проверки (тестирования кода) решим несколько задач
/*
* -----Задача 1:
* Вводится последовательность целых чисел, оканчивающаяся на нулём.
* Сам ноль в последовательность не входит.
* Постройте в этой последовательности дерева.
* Для полученного дерева выведите список всех вершин, 
* имеющих по два ребёнка в порядке возрастания.
* Замечание: если во входных данных есть повторяющие числа, то их нужно 
игнорировать, т.е. все элементы дерева уникальны
* Пример:
* 7 3 2 1 9 5 4 6 8 0
* -----Задача 2:
* Напишите программу, которая будет реализовывать действия в бинарном 
* дереве поиска "вставить" и "найти" (по значению). 
* Программа должна обрабатывать запросы 3х видов:
* 1) ADD n - если указанного числа ещё нет в дереве, вставлять его и 
* выводить слово "DONE", если уже есть - оставлять дерево как было и 
* выводить слово "ALREADY".
* 2) SEARCH - следует выводить слово "YES" 
* (если значение найдено в дереве) или слово "NO" 
* (если не найдено). Дерево при этом не меняется.
* 3) PRINTTREE - выводить всё дерево, обязательно используя алгоритм, 
указанный далее в условии задачи.

Алгоритм для вывода дерева:
template<class T> void print_tree(Node<T> *p, int level)
{
	if (p == nullptr)
		return;
	print_tree(p->left, level + 1);
	for(int i = 0; i < level; i++)
		cout << ".";
	cout << p->data << end;
	print_tree(p->right, level + 1);
}
* Изначальный вызов этой функции - print_tree(root, 0)
* Пример:
ADD 2
ADD 3
ADD 2
SEARCH 2
ADD 5
PRINTTREE
SEARCH 7
*/
int main()
{
	// для задачи 2:
	// Если есть input.txt, то переключаем ввод на него.
	// Раз есть файл с входными данными — значит 
  // предполагается, что есть файл и для 
  // результата, пишем в него.
	// Если output.txt отсутствует — вывод останется в консоли.
	if (freopen("input.txt", "r", stdin)) {
		(void) freopen("output.txt", "w", stdout);
	}
	// Если input.txt не открылся — freopen вернул nullptr,
	// cin по-прежнему связан с консолью, работаем через неё
	BST bst;
	// для 1й задачи

	/*
	while (true)
	{
		int x;
		cin >> x;
		if (x == 0)
		{
			break;
		}
		bst.insert(x);
	}
	bst.print();
	*/

	string command;
	while (cin >> command)
	{
		if (command == "SEARCH")
		{
			int key;
			cin >> key;
			cout << (bst.contains(key) ? "YES\n" : 
      "NO\n");
		}
		else if (command == "ADD")
		{
			int key;
			cin >> key;
			cout << (bst.insert(key) ? "DONE\n" : 
      "ALREADY\n");
		}
		else if (command == "PRINTTREE") 
			bst.print();
		else cout << "Unknown command\n";
	}
}
\end{minted}

\section{Реализация Simple Snippets}
\label{sec:ss}

\subsection{Описание реализации}

Следующая рассматриваемая реализация взята из курса Simple
Snippets~\cite{simplesnippets}. В отличие от наивной версии, она содержит
полный набор операций для работы с деревом: вставку, поиск, удаление, три
вида обхода (прямой, симметричный и обратный), обход по уровням,
вычисление высоты и вывод дерева, повёрнутого на $90^\circ$. Управление
организовано через консольное меню.

Узел представлен классом \texttt{TreeNode} с полями \texttt{value} (ключ),
\texttt{left} и \texttt{right}. Класс \texttt{BST} хранит указатель на корень
\texttt{root} и содержит следующие методы:
\begin{itemize}
  \item \texttt{insertNode} и \texttt{insertRecursive} "--- две реализации
        вставки (итеративная и рекурсивная); повторная вставка существующего
        ключа отклоняется;
  \item \texttt{iterativeSearch} и \texttt{recursiveSearch} "--- две
        реализации поиска;
  \item \texttt{deleteNode} "--- рекурсивное удаление, возвращающее указатель
        на корень поддерева после удаления;
  \item \texttt{minValueNode} "--- поиск минимального узла в поддереве
        (нужен при удалении узла с двумя детьми);
  \item \texttt{printPreorder}, \texttt{printInorder},
        \texttt{printPostorder} "--- три вида обхода;
  \item \texttt{height}, \texttt{printGivenLevel},
        \texttt{printLevelOrderBFS} "--- высота дерева и обход по уровням;
  \item \texttt{print2D} "--- вывод дерева "боком".
\end{itemize}
Код приведён в листинге ниже.

\begin{minted}{c}
#include<iostream>
#define SPACE 10

using namespace std;

class TreeNode {
  public:
    int value;
  TreeNode * left;
  TreeNode * right;

  TreeNode() {
    value = 0;
    left = NULL;
    right = NULL;
  }
  TreeNode(int v) {
    value = v;
    left = NULL;
    right = NULL;
  }
};

class BST {
  public:
    TreeNode * root;
  BST() {
    root = NULL;
  }
  bool isTreeEmpty() {
    if (root == NULL) {
      return true;
    } else {
      return false;
    }
  }

  void insertNode(TreeNode * new_node) {
    if (root == NULL) {
      root = new_node;
      cout << "Value Inserted as root node!" << endl;
    } else {
      TreeNode * temp = root;
      while (temp != NULL) {
        if (new_node -> value == temp -> value) {
          cout << "Value Already exist," <<
            "Insert another value!" << endl;
          return;
        } else if ((new_node -> value < temp -> value) && 
        (temp -> left == NULL)) {
          temp -> left = new_node;
          cout << "Value Inserted to the left!" << endl;
          break;
        } else if (new_node -> value < temp -> value) {
          temp = temp -> left;
        } else if ((new_node -> value > temp -> value) && 
        (temp -> right == NULL)) {
          temp -> right = new_node;
          cout << "Value Inserted to the right!" << endl;
          break;
        } else {
          temp = temp -> right;
        }
      }
    }
  }
  	TreeNode* insertRecursive(TreeNode *r, TreeNode *new_node)
	{
		if(r==NULL)
		{
			r=new_node;
			cout <<"Insertion successful"<<endl;
			return r;
		}
	
		if(new_node->value < r->value)
		{
			r->left = insertRecursive(r->left,new_node);
		}
		else if (new_node->value > r->value)  
		{
			r->right = insertRecursive(r->right,new_node);
		}
	   else
	   {
	     	cout << "No duplicate values allowed!" << endl;
	     	return r; 
		}
		return r;
	}

  void print2D(TreeNode * r, int space) {
    if (r == NULL) // Base case  1
      return;
    space += SPACE; // Increase distance between levels   2
    print2D(r -> right, space); // Process right child first 3 
    cout << endl;
    for (int i = SPACE; i < space; i++) // 5 
      cout << " "; // 5.1  
    cout << r -> value << "\n"; // 6
    print2D(r -> left, space); // Process left child  7
  }

  void printPreorder(TreeNode * r) //(current node, Left, Right) 
  {
    if (r == NULL)
      return;
    /* first print data of node */
    cout << r -> value << " ";
    /* then recur on left sutree */
    printPreorder(r -> left);
    /* now recur on right subtree */
    printPreorder(r -> right);
  }

  void printInorder(TreeNode * r) //(Left, current node, Right)
  {
    if (r == NULL)
      return;
    /* first recur on left child */
    printInorder(r -> left);
    /* then print the data of node */
    cout << r -> value << " ";
    /* now recur on right child */
    printInorder(r -> right);
  }
  void printPostorder(TreeNode * r) //(Left, Right, Root)
  {
    if (r == NULL)
      return;
    // first recur on left subtree 
    printPostorder(r -> left);
    // then recur on right subtree 
    printPostorder(r -> right);
    // now deal with the node 
    cout << r -> value << " ";
  }

  TreeNode * iterativeSearch(int v) {
    if (root == NULL) {
      return root;
    } else {
      TreeNode * temp = root;
      while (temp != NULL) {
        if (v == temp -> value) {
          return temp;
        } else if (v < temp -> value) {
          temp = temp -> left;
        } else {
          temp = temp -> right;
        }
      }
      return NULL;
    }
  }

  TreeNode * recursiveSearch(TreeNode * r, int val) {
    if (r == NULL || r -> value == val)
      return r;

    else if (val < r -> value)
      return recursiveSearch(r -> left, val);

    else
      return recursiveSearch(r -> right, val);
  }

  int height(TreeNode * r) {
    if (r == NULL)
      return -1;
    else {
      /* compute the height of each subtree */
      int lheight = height(r -> left);
      int rheight = height(r -> right);

      /* use the larger one */
      if (lheight > rheight)
        return (lheight + 1);
      else return (rheight + 1);
    }
  }

  /* Print nodes at a given level */
  void printGivenLevel(TreeNode * r, int level) {
    if (r == NULL)
      return;
    else if (level == 0)
      cout << r -> value << " ";
    else // level > 0  
    {
      printGivenLevel(r -> left, level - 1);
      printGivenLevel(r -> right, level - 1);
    }
  }
  void printLevelOrderBFS(TreeNode * r) {
    int h = height(r);
    for (int i = 0; i <= h; i++)
      printGivenLevel(r, i);
  }

  TreeNode * minValueNode(TreeNode * node) {
    TreeNode * current = node;
    /* loop down to find the leftmost leaf */
    while (current -> left != NULL) {
      current = current -> left;
    }
    return current;
  }

  TreeNode * deleteNode(TreeNode * r, int v) {
    // base case 
    if (r == NULL) {
      return NULL;
    }
    // If the key to be deleted is smaller than the root's key, 
    // then it lies in left subtree 
    else if (v < r -> value) {
      r -> left = deleteNode(r -> left, v);
    }
    // If the key to be deleted is greater than the root's key, 
    // then it lies in right subtree 
    else if (v > r -> value) {
      r -> right = deleteNode(r -> right, v);
    }
    // if key is same as root's key, then 
    // This is the node to be deleted 
    else {
      // node with only one child or no child 
      if (r -> left == NULL) {
        TreeNode * temp = r -> right;
        delete r;
        return temp;
      } else if (r -> right == NULL) {
        TreeNode * temp = r -> left;
        delete r;
        return temp;
      } else {
        // node with two children: Get the inorder successor 
        // (smallest in the right subtree) 
        TreeNode * temp = minValueNode(r -> right);
        // Copy the inorder successor's content to this node 
        r -> value = temp -> value;
        // Delete the inorder successor 
        r -> right = deleteNode(r -> right, temp -> value);
        //deleteNode(r->right, temp->value); 
      }
    }
    return r;
  }

};

int main() {
  BST obj;
  int option, val;

  do {
    cout << "What operation do you want to perform? " <<
      " Select Option number. Enter 0 to exit." << endl;
    cout << "1. Insert Node" << endl;
    cout << "2. Search Node" << endl;
    cout << "3. Delete Node" << endl;
    cout << "4. Print/Traversal BST values" << endl;
    cout << "5. Height of Tree" << endl;
    cout << "6. Clear Screen" << endl;
    cout << "0. Exit Program" << endl;

    cin >> option;
    //Node n1;
    TreeNode * new_node = new TreeNode();

    switch (option) {
    case 0:
      break;
    case 1:
      	cout <<"INSERT"<<endl;
	      cout <<"Enter VALUE of TREE NODE to INSERT in BST: ";
	      cin >> val;
	      new_node->value = val;
	      obj.root= obj.insertRecursive(obj.root,new_node);
	      //obj.insertNode(new_node);
	      cout<<endl;
    		break;
      
    case 2:
      cout << "SEARCH" << endl;
      cout << "Enter VALUE of TREE NODE to SEARCH in BST: ";
      cin >> val;
      //new_node = obj.iterativeSearch(val);
      new_node = obj.recursiveSearch(obj.root, val);
      if (new_node != NULL) {
        cout << "Value found" << endl;
      } else {
        cout << "Value NOT found" << endl;
      }
      break;
    case 3:
      cout << "DELETE" << endl;
      cout << "Enter VALUE of TREE NODE to DELETE in BST: ";
      cin >> val;
      new_node = obj.iterativeSearch(val);
      if (new_node != NULL) {
        obj.deleteNode(obj.root, val);
        cout << "Value Deleted" << endl;
      } else {
        cout << "Value NOT found" << endl;
      }
      break;
    case 4:
      cout << "PRINT 2D: " << endl;
      obj.print2D(obj.root, 5);
      cout << endl;
      cout << "Print Level Order BFS: \n";
      obj.printLevelOrderBFS(obj.root);
      cout << endl;
      //	      cout <<"PRE-ORDER: ";
      //	      obj.printPreorder(obj.root);
      //	      cout<<endl;
      //	      cout <<"IN-ORDER: ";
      //	      obj.printInorder(obj.root);
      //	      cout<<endl;
      //	      cout <<"POST-ORDER: ";
      //	      obj.printPostorder(obj.root);
      break;
    case 5:
      cout << "TREE HEIGHT" << endl;
      cout << "Height : " << obj.height(obj.root) << endl;
      break;
    case 6:
      system("cls");
      break;
    default:
      cout << "Enter Proper Option number " << endl;
    }

  } while (option != 0);

  return 0;
}
\end{minted}

\subsection{Выявленные недостатки}

Реализация работает на типичных сценариях, но содержит ошибки и
проектные недостатки. Часть из них проверена запуском программы под
AddressSanitizer и LeakSanitizer (компилятор g++).

\begin{enumerate}
  \item \textbf{Потеря нового корня при удалении (ошибка).} В пункте меню
        «Удалить» результат вызова \texttt{obj.deleteNode(obj.root, val)}
        отбрасывается, хотя функция возвращает указатель на корень поддерева
        после удаления. Если удаляется корень, у которого нет одного из
        потомков, память освобождается оператором \texttt{delete}, а
        \texttt{obj.root} продолжает указывать на неё, и любая следующая
        операция обращается к освобождённой памяти. Проверка: вставка
        ключа и его удаление из дерева из одного узла с последующим выводом
        дают ошибку \texttt{heap-use-after-free}.

  \item \textbf{Утечки памяти.} 
  
  	\textbf{(а)} В начале каждой итерации цикла меню безусловно выполняется \texttt{new TreeNode()}, 
		хотя узел нужен только при вставке. 
		
		\textbf{(б)} В пунктах «Найти» и «Удалить» указатель \texttt{new\_node} перезаписывается результатом 
		поиска, и созданный узел теряется. 
		
		\textbf{(в)} При вставке дубликата новый узел не освобождается. 
		
		\textbf{(г)} У класса \texttt{BST} нет деструктора, поэтому
        при выходе из программы дерево не освобождается. Проверка: после
        вставки трёх ключей и немедленного выхода LeakSanitizer сообщает об
        утечке 96 байт в четырёх блоках "--- три узла дерева и узел,
        созданный на последней итерации меню.

  \item \textbf{Отсутствие проверки ввода.} Значения читаются без проверки
        результата. При вводе нечислового значения поток \texttt{cin}
        переходит в состояние ошибки, последующие операции чтения не
        изменяют переменную \texttt{option}, и программа бесконечно
        повторяет последний выбранный пункт. Проверка: при вводе букв
        вместо ключа ключ 0 вставляется один раз, после чего программа
        бесконечно выводит сообщение о дубликате.

  \item \textbf{Дублирование кода и двойной поиск.} Есть две реализации
        вставки и две "--- поиска, при этом в меню используется только по
        одной (остальные закомментированы или не вызываются). 
		
		При удалении сначала вызывается \texttt{iterativeSearch}, чтобы проверить
        наличие ключа, а затем \texttt{deleteNode} ищет тот же ключ повторно.

  \item \textbf{Смешение логики и ввода-вывода.} Методы вставки сами
        печатают сообщения, поэтому результат операции нельзя использовать
        в программе (например, в тесте) и нельзя повторно использовать класс
        в другой программе. Кроме того, узел создаёт вызывающий код, и
        неясно, кто отвечает за его освобождение.

  \item \textbf{Нарушение инкапсуляции.} Все поля и вспомогательные
        рекурсивные методы объявлены \texttt{public}. Чтобы вставить ключ,
        приходится писать 
		
		\texttt{obj.insertRecursive(obj.root, new\_node)}, и пользователь класса 
		может случайно нарушить свойство дерева, изменив \texttt{root} или указатели узлов.

  \item \textbf{Неэффективный обход по уровням.}
        \texttt{printLevelOrderBFS} для каждого уровня заново вызывает
        \texttt{printGivenLevel}, который обходит дерево с корня. Суммарное
        время составляет $O(n\cdot h)$, а для вырожденного дерева
        $O(n^2)$. Название BFS вводит в заблуждение: очереди нет, по сути
        выполняется повторяющийся поиск в глубину.

  \item \textbf{Устаревший и непереносимый стиль.} Используются
        \texttt{NULL} вместо \texttt{nullptr}, макрос
        \texttt{\#define SPACE}, очистка экрана через
        \texttt{system("cls")} (работает только в Windows); методы, не
        изменяющие дерево, не помечены \texttt{const}; копирование объекта
        не запрещено.

  \item \textbf{Глубина рекурсии.} Рекурсивная вставка, удаление, обходы и
        вычисление высоты используют стек, глубина которого равна высоте
        дерева. Для вырожденного дерева из очень большого числа узлов
        это грозит переполнением стека. Недостаток свойственен всем
        рекурсивным реализациям и обсуждается в разделе~\ref{sec:analysis}.
\end{enumerate}

\subsection{Способы исправления}

Для каждого недостатка ниже указан способ его устранения; реализация с
этими исправлениями описана в разделе~\ref{sec:good}.

\begin{longtable}{|c|p{0.28\textwidth}|p{0.54\textwidth}|}
  \caption{Недостатки реализации Simple Snippets и их исправление}
  \label{tab:fixes} \\
  \hline
  № & Недостаток & Исправление \\ \hline
  \endfirsthead
  \hline
  № & Недостаток & Исправление \\ \hline
  \endhead
  1 & Потеря нового корня при удалении &
      Передавать указатель по ссылке (\texttt{Node*\&}): присваивание
      \texttt{node = ...} изменяет указатель в родителе или сам
      \texttt{root}, возвращать значение не нужно. \\ \hline
  2 & Утечки памяти &
      Узлы создаёт и удаляет только сам класс: \texttt{insert(int key)}
      создаёт узел внутри, лишних \texttt{new} в меню нет; деструктор
      освобождает дерево обратным обходом; копирование запрещено
      (\texttt{= delete}). \\ \hline
  3 & Нет проверки ввода &
      Функция \texttt{readInt} проверяет результат \texttt{cin >{}> value};
      при ошибке сбрасывает флаги (\texttt{cin.clear()}), пропускает
      остаток строки (\texttt{cin.ignore}) и выводит сообщение. \\ \hline
  4 & Дублирование кода, двойной поиск &
      Оставлен один вариант каждой операции; \texttt{remove} находит и
      удаляет узел за один спуск и возвращает \texttt{bool}. \\ \hline
  5 & Логика смешана с выводом &
      Методы возвращают \texttt{bool} (выполнено или нет), печатает
      только \texttt{main}; в метод вставки передаётся ключ, а не узел. \\ \hline
  6 & Нарушена инкапсуляция &
      \texttt{root}, \texttt{Node} и рекурсивные помощники
      \texttt{private}; публичный интерфейс состоит из десяти методов. \\ \hline
  7 & Обход по уровням за $O(n\cdot h)$ &
      Использовать очередь \texttt{std::queue}: каждый узел помещается в
      неё и извлекается ровно один раз, время $O(n)$. \\ \hline
  8 & Устаревший стиль &
      \texttt{nullptr} вместо \texttt{NULL}; макрос заменён выражением;
      \texttt{system("cls")} удалён; методы, не меняющие дерево, помечены
      \texttt{const}. \\ \hline
  9 & Глубина рекурсии &
      Не устранён: поиск сделан итеративным, остальные операции остаются
      рекурсивными (см. подраздел~\ref{sec:limits}). \\ \hline
\end{longtable}


% =====================================================================
\section{Хорошая, правильная реализация}
\label{sec:good}

Улучшенная реализация сохраняет возможности версии Simple Snippets
(вставка, поиск, удаление, три вида обхода, обход по уровням, высота,
вывод дерева, меню), но устраняет недостатки, перечисленные в
таблице~\ref{tab:fixes}. Весь код приведён в листинге ниже, далее
разобрано, как он работает.

\begin{minted}{c}
// Двоичное дерево поиска (BST) на C++17.
// Без windows.h: код работает на любой ОС.

#include <iostream>
#include <limits>
#include <queue>
#include <string>
using namespace std;

class BST {
public:
    BST() = default;
    ~BST() { clear(root); } // освобождаем всё дерево

    // Копировать дерево запрещаем: иначе два объекта 
	// владели бы одними узлами
    BST(const BST&) = delete;
    BST& operator=(const BST&) = delete;

    // ---- Основные операции. Возвращают true/false, 
	// ничего не печатают ----
    bool insert(int key) { return insertHelper(root, key); }
    bool remove(int key) { return removeHelper(root, key); }
    bool contains(int key) const 
	{ 
		return findHelper(root, key) != nullptr; 
	}
    bool empty() const { return root == nullptr; }

    // ---- Три обхода ----
    void preorder() const 
	{ 
		preorderHelper(root);  
		cout << '\n'; 
	}   // узел, левое, правое
    void inorder() const 
	{ 
		inorderHelper(root);   
		cout << '\n'; 
	}   // левое, узел, правое
    void postorder() const 
	{ 
		postorderHelper(root); 
		cout << '\n'; 
	}   // левое, правое, узел

    // ---- Обход по уровням (через очередь, за O(n)) ----
    void levelOrder() const {
        if (!root) { cout << '\n'; return; }
        queue<Node*> q;
        q.push(root);
        while (!q.empty()) {
            Node* cur = q.front();
            q.pop();
            cout << cur->key << ' ';
            if (cur->left)  q.push(cur->left);
            if (cur->right) q.push(cur->right);
        }
        cout << '\n';
    }
    // Высота в рёбрах: пустое дерево -1, одиночный узел 0
    int height() const { return heightHelper(root); }
    // Дерево, повёрнутое на 90 градусов: корень слева, 
	// правые дети сверху
    void print2D() const {
        if (!root) { cout << "(пусто)\n"; return; }
        print2DHelper(root, 0);
    }

private:
    struct Node {
        int key;
        Node* left = nullptr;
        Node* right = nullptr;
        explicit Node(int k) : key(k) {}
    };

    Node* root = nullptr;

    // ---- Вставка ----
    // Node*& значит "ссылка на сам указатель". 
    // Присвоив node = new Node(...),
    // мы меняем тот указатель, который лежит в родителе (или root).
    bool insertHelper(Node*& node, int key) {
        if (!node) {	// нашли пустое место
            node = new Node(key);
            return true;
        }
        if (key < node->key)
			return insertHelper(node->left, key);
        if (key > node->key) 
			return insertHelper(node->right, key);
        return false; // такой ключ уже есть
    }

    // ---- Поиск ----
    Node* findHelper(Node* node, int key) const {
        while (node && node->key != key)
            node = (key < node->key) ? node->left : node->right;
        return node;	// nullptr, если не нашли
    }

    // ---- Удаление ----
    bool removeHelper(Node*& node, int key) {
        if (!node) return false;	// не нашли
        if (key < node->key) return removeHelper(node->left, key);
        if (key > node->key) return removeHelper(node->right, key);

        // Нашли удаляемый узел.
		// нет левого: заменяем правым (или nullptr)
        if (!node->left) {
            Node* old = node;
            node = node->right;
            delete old;
        }
		// нет правого: заменяем левым
        else if (!node->right) {
            Node* old = node;
            node = node->left;
            delete old;
        }
        else {// два ребёнка
			// преемник = минимум правого поддерева
            Node* succ = node->right;
            while (succ->left) succ = succ->left;
            node->key = succ->key;	// копируем его ключ сюда
			// и удаляем преемника (у него нет левого)
            removeHelper(node->right, succ->key);
        }
        return true;
    }

    // ---- Обходы ----
    void preorderHelper(const Node* node) const {
        if (!node) return;
        cout << node->key << ' ';
        preorderHelper(node->left);
        preorderHelper(node->right);
    }
    void inorderHelper(const Node* node) const {
        if (!node) return;
        inorderHelper(node->left);
        cout << node->key << ' ';
        inorderHelper(node->right);
    }
    void postorderHelper(const Node* node) const {
        if (!node) return;
        postorderHelper(node->left);
        postorderHelper(node->right);
        cout << node->key << ' ';
    }

    int heightHelper(const Node* node) const {
        if (!node) return -1;
        int l = heightHelper(node->left);
        int r = heightHelper(node->right);
        return 1 + (l > r ? l : r);
    }

    void print2DHelper(const Node* node, int depth) const {
        if (!node) return;
		// сначала правое: оно окажется сверху
        print2DHelper(node->right, depth + 1);
        cout << string(depth * 4, ' ') << node->key << '\n';
        print2DHelper(node->left, depth + 1);
    }

    // Удаление всего поддерева обратным обходом: 
	// сначала дети, потом сам узел
    void clear(Node* node) {
        if (!node) return;
        clear(node->left);
        clear(node->right);
        delete node;
    }
};

// Читает целое число; при неверном вводе очищает поток 
// и возвращает false
bool readInt(const char* prompt, int& value) {
    cout << prompt;
    if (cin >> value) return true;
    if (cin.eof()) exit(0);	// конец ввода: выходим
    cin.clear();
    cin.ignore(numeric_limits<streamsize>::max(), '\n');
    cout << "Нужно ввести целое число.\n";
    return false;
}

int main() 
{
    setlocale(LC_ALL, "Russian");
    BST tree;
    int option = -1;

    while (option != 0) {
        cout << "\n1. Вставить  2. Найти  3. Удалить  
		4. Обходы и вывод  5. Высота  0. Выход\n";
        if (!readInt("> ", option)) { option = -1; continue; }

        int val;
        switch (option) {
        case 0:
            break;
        case 1:
            if (readInt("Ключ: ", val))
                cout << (tree.insert(val) ? "Вставлено\n" : 
				"Такой ключ уже есть\n");
            break;
        case 2:
            if (readInt("Ключ: ", val))
                cout << (tree.contains(val) ? "Найден\n" : 
				"Не найден\n");
            break;
        case 3:
            if (readInt("Ключ: ", val))
                cout << (tree.remove(val) ? "Удалён\n" : "Не найден\n");
            break;
        case 4:
            cout << "Прямой:      "; tree.preorder();
            cout << "Симметричный: "; tree.inorder();
            cout << "Обратный:    "; tree.postorder();
            cout << "По уровням:  "; tree.levelOrder();
            cout << "Дерево:\n";    tree.print2D();
            break;
        case 5:
            cout << "Высота: " << tree.height() << '\n';
            break;
        default:
            cout << "Нет такого пункта\n";
        }
    }
    return 0;
}
\end{minted}

\subsection{Устройство класса}

Класс \texttt{BST} делится на открытый интерфейс и закрытую часть.
Снаружи видны только операции над деревом, а узлы, корень и рекурсивные
помощники скрыты.

\begin{longtable}{|p{0.25\textwidth}|p{0.5\textwidth}|p{0.15\textwidth}|}
  \caption{Открытый интерфейс класса \texttt{BST}}
  \label{tab:interface} \\
  \hline
  Метод & Назначение & Результат \\ \hline
  \endfirsthead
  \hline
  Метод & Назначение & Результат \\ \hline
  \endhead
  \texttt{insert(key)}   & вставка ключа; повторный ключ отклоняется & \texttt{bool} \\ \hline
  \texttt{remove(key)}   & удаление узла с заданным ключом & \texttt{bool} \\ \hline
  \texttt{contains(key)} & проверка наличия ключа & \texttt{bool} \\ \hline
  \texttt{empty()}       & проверка, что дерево пусто & \texttt{bool} \\ \hline
  \texttt{preorder()}, \texttt{inorder()}, \texttt{postorder()} &
      прямой, симметричный и обратный обходы с печатью ключей & \texttt{void} \\ \hline
  \texttt{levelOrder()}  & обход по уровням (в ширину) & \texttt{void} \\ \hline
  \texttt{height()}      & высота дерева в рёбрах (пустое дерево $-1$) & \texttt{int} \\ \hline
  \texttt{print2D()}     & вывод дерева, повёрнутого на $90^\circ$ & \texttt{void} \\ \hline
\end{longtable}

Узел \texttt{Node} описан внутри класса как закрытая структура с ключом и
двумя указателями на потомков. Корень \texttt{root} также закрыт, поэтому
все узлы создаются и удаляются только методами самого класса. Это
исключает вопрос о владении памятью, который в версии Simple Snippets
оставался нерешённым.

\subsection{Вставка}

Вставка выполнена рекурсивной функцией, принимающей указатель на узел
\emph{по ссылке}:

\begin{minted}[fontsize=\small]{cpp}
bool insertHelper(Node*& node, int key) {
    if (!node) {                              // нашли пустое место
        node = new Node(key);
        return true;
    }
    if (key < node->key) return insertHelper(node->left, key);
    if (key > node->key) return insertHelper(node->right, key);
    return false;                             // такой ключ уже есть
}
\end{minted}

Запись \texttt{Node*\&} означает ссылку на сам указатель, а не на узел,
на который он указывает. Поэтому присваивание \texttt{node = new Node(key)}
изменяет то поле, которое было передано: \texttt{root} либо
\texttt{left}/\texttt{right} родителя. Отдельно хранить указатель на
родителя и возвращать новый корень поддерева не требуется.

Пример: вставка ключа 10 в дерево с ключами 7, 3, 2, 1, 9, 5, 4, 6, 8.
Так как $10 > 7$, вызов переходит к правому потомку; затем $10 > 9$, и
вызов получает ссылку на правое поле узла 9, которое равно
\texttt{nullptr}. Там создаётся узел 10, и он оказывается правым потомком
девятки. Вставка уже имеющегося ключа доходит до последней строки и
возвращает \texttt{false}, дерево не меняется.

\subsection{Поиск}

Поиск выполнен итеративно и не использует стек вызовов:

\begin{minted}[fontsize=\small]{cpp}
Node* findHelper(Node* node, int key) const {
    while (node && node->key != key)
        node = (key < node->key) ? node->left : node->right;
    return node;                              // nullptr, если не нашли
}
\end{minted}

Цикл идёт вниз, пока не встретит искомый ключ или пустой указатель.
Метод \texttt{contains} возвращает \texttt{true}, если результат не равен
\texttt{nullptr}.

\subsection{Удаление}

Удаление использует ту же идею ссылки на указатель и рассматривает три
случая~\cite{kormen}:

\begin{minted}[fontsize=\small]{cpp}
bool removeHelper(Node*& node, int key) {
    if (!node) return false;                  // не нашли
    if (key < node->key) return removeHelper(node->left, key);
    if (key > node->key) return removeHelper(node->right, key);

    if (!node->left) {                        // нет левого ребёнка
        Node* old = node;
        node = node->right;
        delete old;
    } else if (!node->right) {                // нет правого ребёнка
        Node* old = node;
        node = node->left;
        delete old;
    } else {                                  // два ребёнка
        Node* succ = node->right;             // преемник: минимум справа
        while (succ->left) succ = succ->left;
        node->key = succ->key;                // копируем его ключ
        removeHelper(node->right, succ->key); // удаляем преемника
    }
    return true;
}
\end{minted}

\begin{itemize}
  \item \textbf{Нет левого ребёнка} (в том числе лист). Узел заменяется
        своим правым потомком, а если его нет, то \texttt{nullptr}. Так как
        \texttt{node} "--- ссылка на указатель родителя, замена выполняется
        одним присваиванием.
  \item \textbf{Нет правого ребёнка.} Симметрично, узел заменяется левым
        потомком. Пример: узел 9 в дереве выше имеет единственного потомка 8,
        который встаёт на его место.
  \item \textbf{Два ребёнка.} Нельзя просто убрать узел, поэтому на его
        место ставится преемник "--- следующий по величине ключ, то есть
        минимум правого поддерева. Он достигается шагом вправо и спуском
        влево до конца. Ключ преемника копируется в удаляемый узел, после
        чего сам преемник удаляется из правого поддерева; у него нет левого
        ребёнка, поэтому это один из двух предыдущих случаев.
\end{itemize}

Пример. При удалении ключа 5 (дети 4 и 6) преемником становится 6: его
ключ записывается в узел, который раньше хранил 5, а старый узел 6
удаляется. При удалении корня 7 преемником становится 8 "--- минимум
правого поддерева, корнем которого служит 9. Так как удаление корня тоже
выполняется через ссылку на \texttt{root}, дерево остаётся корректным,
в отличие от версии Simple Snippets.

\subsection{Обходы дерева}

Три вида обхода отличаются только положением печати ключа относительно
двух рекурсивных вызовов:
\begin{itemize}
  \item прямой (\texttt{preorder}): узел, левое поддерево, правое поддерево;
  \item симметричный (\texttt{inorder}): левое поддерево, узел, правое
        поддерево; ключи выводятся по возрастанию, что используется как
        проверка корректности дерева;
  \item обратный (\texttt{postorder}): левое поддерево, правое поддерево,
        узел.
\end{itemize}
Для дерева с ключами 7, 3, 2, 1, 9, 5, 4, 6, 8 они выводят
\texttt{7 3 2 1 5 4 6 9 8}, \texttt{1 2 3 4 5 6 7 8 9} и
\texttt{1 2 4 6 5 3 8 9 7} соответственно. Время каждого обхода равно
$\Theta(n)$, как показано в теореме о времени центрированного обхода из раздела «Краткая теория».

\subsection{Обход по уровням и высота}

Обход по уровням выполнен с помощью очереди:

\begin{minted}[fontsize=\small]{cpp}
queue<Node*> q;
q.push(root);
while (!q.empty()) {
    Node* cur = q.front();
    q.pop();
    cout << cur->key << ' ';
    if (cur->left)  q.push(cur->left);
    if (cur->right) q.push(cur->right);
}
\end{minted}

Сначала в очередь помещается корень. На каждом шаге из очереди
извлекается узел, печатается его ключ, а его потомки добавляются в конец.
Так узлы выводятся уровень за уровнем, слева направо. Каждый узел
помещается в очередь и извлекается ровно один раз, поэтому время равно
$O(n)$, в отличие от $O(n \cdot h)$ в версии Simple Snippets. Для того же
дерева результат: \texttt{7 3 9 2 5 8 1 4 6}.

Высота считается рекурсивно как единица плюс большая из высот поддеревьев
и измеряется в рёбрах: пустое дерево имеет высоту $-1$, дерево из одного
узла "--- 0. Такое соглашение совпадает с принятым в~\cite{kormen}.

\subsection{Вывод дерева}

Метод \texttt{print2D} выводит дерево, повёрнутое на $90^\circ$:
корень находится слева, правые потомки "--- выше, левые "--- ниже. Для этого
выполняется обход «правое поддерево, узел, левое поддерево», а перед
каждым ключом печатается отступ в $4 \cdot d$ пробелов, где $d$ "--- глубина
узла. Благодаря этому форма дерева видна сразу, без \texttt{windows.h}, и
программа работает на любой операционной системе.

\subsection{Освобождение памяти}

Деструктор вызывает функцию \texttt{clear}, которая обходит дерево в
обратном порядке: сначала удаляются оба поддерева, затем сам узел. Порядок
важен, поскольку после \texttt{delete} указатели на потомков уже
недоступны. Копирование объекта запрещено (\texttt{BST(const BST\&) = delete}
и такой же оператор присваивания), иначе два объекта владели бы одними и
теми же узлами и попытались бы удалить их дважды.

\subsection{Ввод и меню}

Функция \texttt{readInt} читает целое число и проверяет результат:
при неверном вводе она сбрасывает состояние потока, пропускает остаток
строки, сообщает об ошибке и возвращает \texttt{false}, а при конце ввода
завершает программу. В отличие от версии Simple Snippets, меню продолжает
работать после неверного ввода. Операции вызываются через
\texttt{switch}; сообщения печатает только \texttt{main}, а методы класса
возвращают \texttt{bool}.

\subsection{Оставшиеся ограничения}
\label{sec:limits}

\begin{itemize}
  \item Вставка, удаление, обходы, высота и вывод остаются рекурсивными,
        поэтому глубина стека равна высоте дерева. Поиск итеративный.
        Рекурсию можно заменить циклом или явным стеком, но это усложнило
        бы код, и в данной работе не делается.
  \item Дерево не балансируется: форма зависит от порядка вставки
        (раздел~\ref{sec:analysis}).
  \item Ключи имеют тип \texttt{int}; для произвольного типа класс можно
        сделать шаблоном.
  \item Высота не хранится в узлах и вычисляется заново при каждом
        вызове за $\Theta(n)$.
\end{itemize}


% =====================================================================
\section{Демонстрация работы}
\label{sec:demo}

Ниже приведены результаты работы улучшенной реализации. Пункты меню в
листингах опущены, показан только результат выбранных операций.

\textbf{1. Построение дерева.} В пустое дерево вставляются ключи
7, 3, 2, 1, 9, 5, 4, 6, 8, затем выбирается пункт «Обходы и вывод»:

\begin{Verbatim}[fontsize=\small]
Прямой:       7 3 2 1 5 4 6 9 8
Симметричный: 1 2 3 4 5 6 7 8 9
Обратный:     1 2 4 6 5 3 8 9 7
По уровням:   7 3 9 2 5 8 1 4 6
Дерево:
    9
        8
7
            6
        5
            4
    3
        2
            1
\end{Verbatim}

Симметричный обход выдаёт ключи по возрастанию, значит, свойство дерева
поиска выполнено. Высота дерева, как показывает пункт «Высота», равна 3.

\textbf{2. Повторная вставка и поиск.} Попытка вставить ключ 5 повторно
отклоняется, поиск находит ключ 6 и не находит ключ 10:

\begin{Verbatim}[fontsize=\small]
Вставка 5:  Такой ключ уже есть
Поиск 6:    Найден
Поиск 10:   Не найден
\end{Verbatim}

\textbf{3. Удаление узла с двумя детьми (ключ 5).} Преемником становится
ключ 6, он занимает место удалённого узла:

\begin{Verbatim}[fontsize=\small]
Удалён
Прямой:       7 3 2 1 6 4 9 8
Симметричный: 1 2 3 4 6 7 8 9
Обратный:     1 2 4 6 3 8 9 7
По уровням:   7 3 9 2 6 8 1 4
Дерево:
    9
        8
7
        6
            4
    3
        2
            1
\end{Verbatim}

\textbf{4. Удаление корня (ключ 7).} Новым корнем становится преемник 8:

\begin{Verbatim}[fontsize=\small]
Удалён
Прямой:       8 3 2 1 6 4 9
Симметричный: 1 2 3 4 6 8 9
Обратный:     1 2 4 6 3 9 8
По уровням:   8 3 9 2 6 1 4
Дерево:
    9
8
        6
            4
    3
        2
            1
\end{Verbatim}

\textbf{5. Удаление листа и несуществующего ключа.} Ключ 1 (лист)
удаляется, ключ 100 не найден, после чего высота остаётся равной 3, а
симметричный обход по-прежнему упорядочен:

\begin{Verbatim}[fontsize=\small]
Удаление 1:    Удалён
Удаление 100:  Не найден
Высота:        3
Симметричный:  2 3 4 6 8 9
\end{Verbatim}

\textbf{6. Неверный ввод.} Если вместо ключа ввести буквы, программа
сообщает об ошибке и продолжает работу (в версии Simple Snippets она
зацикливалась):

\begin{Verbatim}[fontsize=\small]
Ключ: abc
Нужно ввести целое число.
\end{Verbatim}

\textbf{7. Вырождение дерева.} Если вставлять ключи 1, 2, 3, 4, 5 в
порядке возрастания, дерево превращается в цепочку, и его высота равна 4
(на единицу меньше числа узлов):

\begin{Verbatim}[fontsize=\small]
                5
            4
        3
    2
1
\end{Verbatim}

Последний пример объясняет, зачем нужны самобалансирующиеся деревья,
которым посвящены следующие работы.


% =====================================================================
\section{Анализ сложности и расхода памяти}
\label{sec:analysis}

\subsection{Время работы}

Пусть $n$ "--- число узлов, $h$ "--- высота дерева. Поиск, вставка и
удаление спускаются от корня вниз и на каждом уровне выполняют
$O(1)$ действий, поэтому их время пропорционально высоте~\cite{kormen}.
Удаление узла с двумя детьми включает до трёх спусков "--- к удаляемому узлу,
к преемнику и повторный спуск при удалении преемника, "--- так что его
время также равно $O(h)$.

\begin{longtable}{|p{0.38\textwidth}|c|c|c|}
  \caption{Временная сложность операций улучшенной реализации}
  \label{tab:time} \\
  \hline
  Операция & Лучший случай & Средний случай & Худший случай \\ \hline
  \endfirsthead
  \hline
  Операция & Лучший случай & Средний случай & Худший случай \\ \hline
  \endhead
  \texttt{insert}, \texttt{contains}, \texttt{remove} &
      $O(1)$ & $O(\lg n)$ & $O(n)$ \\ \hline
  \texttt{preorder}, \texttt{inorder}, \texttt{postorder} &
      $\Theta(n)$ & $\Theta(n)$ & $\Theta(n)$ \\ \hline
  \texttt{levelOrder} & $\Theta(n)$ & $\Theta(n)$ & $\Theta(n)$ \\ \hline
  \texttt{height} & $\Theta(n)$ & $\Theta(n)$ & $\Theta(n)$ \\ \hline
  \texttt{print2D} (число вызовов) & $\Theta(n)$ & $\Theta(n)$ & $\Theta(n)$ \\ \hline
  деструктор & $\Theta(n)$ & $\Theta(n)$ & $\Theta(n)$ \\ \hline
\end{longtable}

\textbf{Лучший случай} для вставки, поиска и удаления: искомый ключ
находится в корне либо дерево пусто, тогда достаточно одного шага.

\textbf{Средний случай.} Для дерева, построенного вставкой ключей в
случайном порядке, математическое ожидание высоты равно
$O(\lg n)$~\cite{kormen}, поэтому основные операции выполняются в
среднем за $O(\lg n)$, а построение всего дерева "--- за $O(n \lg n)$.

\textbf{Худший случай.} Если ключи вставляются в отсортированном порядке,
каждый новый ключ становится потомком предыдущего, дерево превращается в
цепочку (пример~7 из раздела~\ref{sec:demo}), $h = n-1$, и операции
выполняются за $O(n)$. Построение всего дерева в этом случае занимает
$\sum_{i=1}^{n} i = \Theta(n^2)$.

Все обходы, вычисление высоты и очистка посещают каждый узел ровно один
раз, поэтому их время равно $\Theta(n)$ независимо от формы дерева
(теорема из раздела «Краткая теория»). Метод \texttt{print2D} также делает $\Theta(n)$
вызовов, но объём выводимого текста растёт с глубиной узлов, так как
перед ключом печатается $4d$ пробелов; для цепочки из $n$ узлов он
составляет $\Theta(n^2)$ символов.

Обход по уровням в реализации Simple Snippets работает за $O(n \cdot h)$,
то есть до $O(n^2)$ в худшем случае, а в улучшенной "--- за $\Theta(n)$.

\subsection{Расход памяти}

\begin{itemize}
  \item \textbf{Узлы.} Каждый узел хранит ключ \texttt{int} (4 байта) и
        два указателя. На 64-разрядной платформе \texttt{sizeof(Node)}
        равен 24 байтам (4 байта ключа, 4 байта выравнивания и
        $2 \cdot 8$ байт указателей); к этому добавляются накладные расходы
        аллокатора. Всего потребуется $\Theta(n)$ памяти, объект класса
        \texttt{BST} дополнительно содержит один указатель.
  \item \textbf{Рекурсия.} Вставка, удаление, обходы, вычисление высоты,
        вывод и очистка используют стек вызовов глубиной $h+1$, то есть
        $O(h)$ дополнительной памяти: $O(\lg n)$ в среднем и $O(n)$ в
        худшем случае. Размер стека ограничен (порядка мегабайт, в
        зависимости от системы), поэтому на цепочке из очень большого
        числа узлов рекурсивные операции завершатся аварийно. Поиск
        итеративный и требует $O(1)$ памяти.
  \item \textbf{Обход по уровням.} Очередь одновременно содержит узлы не
        более чем двух соседних уровней, поэтому её размер $O(w)$, где $w$ "---
        максимальная ширина дерева, и в худшем случае $O(n)$.
  \item \textbf{Версия Simple Snippets.} Память под узлы расходуется так
        же, но каждая итерация меню, кроме успешной вставки, теряет
        24 байта, а при выходе не освобождается всё дерево, то есть
        утечка не ограничена сверху.
\end{itemize}

\subsection{Сравнение реализаций}

\begin{longtable}{|p{0.34\textwidth}|c|c|c|}
  \caption{Сравнение трёх реализаций}
  \label{tab:compare} \\
  \hline
  Возможность & Наивная & Simple Snippets & Улучшенная \\ \hline
  \endfirsthead
  \hline
  Возможность & Наивная & Simple Snippets & Улучшенная \\ \hline
  \endhead
  Вставка и поиск                & + & + & + \\ \hline
  Удаление                       & "--- & + (с ошибкой) & + \\ \hline
  Прямой и обратный обходы       & "--- & + & + \\ \hline
  Обход по уровням за $O(n)$     & "--- & "--- ($O(nh)$) & + \\ \hline
  Освобождение памяти            & + & "--- & + \\ \hline
  Проверка ввода                 & "--- & "--- & + \\ \hline
  Защита от копирования          & + & "--- & + \\ \hline
  Балансировка                   & "--- & "--- & "--- \\ \hline
\end{longtable}

\subsection{Вывод}

Время основных операций определяется высотой дерева, а высота зависит от
порядка вставки ключей: от $\lg n$ при случайном порядке до $n$ при
упорядоченном. Улучшенная реализация устраняет ошибки и утечки памяти
версии Simple Snippets и снижает сложность обхода по уровням, но
гарантировать логарифмическую высоту не может. Эту задачу решают
красно-чёрные и АВЛ-деревья.

\conclusion

В ходе работы изучено двоичное дерево поиска: его определение и свойство,
основные операции (поиск, вставка, удаление) и три вида обхода. Показано,
что время основных операций пропорционально высоте дерева, а высота зависит
от порядка вставки ключей.
 
Реализованы и сопоставлены три версии на языке C++: наивная (вставка,
поиск, вывод), версия Simple Snippets (удаление, три вида обхода, обход по
уровням, высота) и улучшенная. В версии Simple Snippets найдены и проверены
запуском недостатки: потеря корня при удалении (обращение к освобождённой
памяти), утечки памяти, зацикливание при неверном вводе, обход по уровням за
$O(n \cdot h)$ и другие. Улучшенная версия их устраняет: указатель передаётся
по ссылке, память создаётся и освобождается только классом, ввод
проверяется, обход по уровням выполняется с помощью очереди за $O(n)$.
 
Оценка сложности показала, что поиск, вставка и удаление выполняются за
$O(h)$, то есть за $O(\lg n)$ в среднем и $O(n)$ в худшем случае, а обходы
и вычисление высоты "--- за $\Theta(n)$. Память составляет $\Theta(n)$ на
узлы (24 байта на узел при 64-разрядной архитектуре) и $O(h)$ на стек
вызовов. Пример с упорядоченной вставкой ключей показал вырождение дерева в
цепочку, поэтому в следующих работах будут рассмотрены самобалансирующиеся
деревья: красно-чёрное и АВЛ-дерево.

% Отобразить все источники. Даже те, на которые нет ссылок.
%\nocite{*}

%\begin{enumerate}
%	\item \href{https://neerc.ifmo.ru/wiki/index.php?title=%D0%94%D0%B5%D1%80%D0%B5%D0%B2%D0%BE_%D0%BF%D0%BE%D0%B8%D1%81%D0%BA%D0%B0,_%D0%BD%D0%B0%D0%B8%D0%B2%D0%BD%D0%B0%D1%8F_%D1%80%D0%B5%D0%B0%D0%BB%D0%B8%D0%B7%D0%B0%D1%86%D0%B8%D1%8F\#.D0.A0.D0.B5.D0.B0.D0.BB.D0.B8.D0.B7.D0.B0%D1%86%D0%B8%D1%8F_.D0.B1.D0.B5.D0.B7_.D0.B8.D1.81%D0%BF%D0%BE%D0%BB%D1%8C%D0%B7%D0%BE%D0%B2%D0%B0%D0%BD%D0%B8%D1%8F_.D0.B8.D0.BD.D1.84.D0.BE.D1.80.D0.BC.%D0%B0%D1%86%D0%B8%D0%B8_.%D0%BE_.%D1%80%D0%BE%D0%B4%D0%B8%D1%82%D0%B5%D0%BB%D0%B5}{Конспект ИТМО по теме "Бинарное дерево поиска, наивная реализация"}
%	\item Simple Snippets (англоязычный автор, соответсвенно все материалы на английском):
%	\begin{itemize}
%		\item \href{https://simplesnippets.in/courses/data-structures-algorithms/}{Структуры данных и алгоритмы
%		Эта серия уроков предназначена для всех, кто хочет изучить тему "Структуры данных и алгоритмы", начав с самых основ.}
%		%\item \href{https://web.archive.org/web/20231227033848/https://simplesnippets.tech/what-is-binary-search-tree-bst-with-full-code-dsa/}{Что такое структура данных "Бинарное дерево поиска" (BST)? Все операции с BST и полный код | Алгоритмы и структуры данных.}
%	\end{itemize}
%	\item \href{https://www.youtube.com/watch?v=0BUX_PotA4c}{Alek OS, "КАК РАБОТАЮТ ДЕРЕВЬЯ | СТРУКТУРЫ ДАННЫХ"}
%\end{enumerate}

\bibliographystyle{ugost2003}
\bibliography{thesis}

% Окончание основного документа и начало приложений Каждая последующая секция
% документа будет являться приложением
\appendix

\end{document}
